\clearpage
\SourceLeaf{614}
\ArithmeticHead{585}{R}{INHALT.}
\begingroup
\newcommand{\ContentsRow}[5]{\par\noindent\hangindent=#1\hspace*{#1}#2\hspace{.4em}\mbox{\footnotesize #3}\hfill\mbox{#4\quad #5}\par}
\ContentsRow{0pt}{59.\quad Integralrechnung, 1797 März 2}{.\,.\,.\,.\,.\,.\,.\,.\,.\,.\,.\,.\,.\,.\,.\,.\,.\,.\,.\,.\,.\,.}{Seite}{514}
\ContentsRow{0pt}{60.\quad Teilung der Lemniskate, 1797 März 19}{.\,.\,.\,.\,.\,.\,.\,.\,.\,.\,.\,.\,.\,.\,.\,.\,.\,.\,.}{—}{515}
\ContentsRow{0pt}{61.\quad Eine Potenzreihe aus der Lehre der lemniskatischen Funktionen, 1797 März}{.\,.\,.\,.\,.}{—}{515}
\ContentsRow{0pt}{62.\quad Fünfteilung der Lemniskate, 1797 März 21}{.\,.\,.\,.\,.\,.\,.\,.\,.\,.\,.\,.\,.\,.\,.\,.\,.}{—}{517}
\ContentsRow{0pt}{63.\quad Über die lemniskatischen Funktionen, 1797 März 29}{.\,.\,.\,.\,.\,.\,.\,.\,.\,.\,.\,.\,.\,.}{—}{517}
\ContentsRow{0pt}{64.\quad Allgemeine Restcharaktere, 1797 Juni 17}{.\,.\,.\,.\,.\,.\,.\,.\,.\,.\,.\,.\,.\,.\,.\,.\,.\,.}{—}{518}
\ContentsRow{0pt}{65.\quad Zur Lehre von den regelmäßigen Vielecken, 1797 Juni 17}{.\,.\,.\,.\,.\,.\,.\,.\,.\,.\,.\,.}{—}{518}
\ContentsRow{0pt}{66.\quad Algebraische Lösung der Kreisteilungsgleichung, 1797 Juli}{.\,.\,.\,.\,.\,.\,.\,.\,.\,.\,.\,.}{—}{518}
\ContentsRow{0pt}{67.\quad Vergl.\ Nr.\ 39, 1797 Juli 20}{.\,.\,.\,.\,.\,.\,.\,.\,.\,.\,.\,.\,.\,.\,.\,.\,.\,.\,.\,.\,.\,.\,.}{—}{519}
\ContentsRow{0pt}{68.\quad Binomische Kongruenz, 1797 Juli 21}{.\,.\,.\,.\,.\,.\,.\,.\,.\,.\,.\,.\,.\,.\,.\,.\,.\,.\,.\,.}{—}{519}
\ContentsRow{0pt}{69.\quad Teilbarkeit ganzer Funktionen, 1797 Juli 23}{.\,.\,.\,.\,.\,.\,.\,.\,.\,.\,.\,.\,.\,.\,.\,.\,.}{—}{519}
\ContentsRow{0pt}{70.\quad Kreisteilungszahlen, 1797 Juli}{.\,.\,.\,.\,.\,.\,.\,.\,.\,.\,.\,.\,.\,.\,.\,.\,.\,.\,.\,.\,.\,.}{—}{520}
\ContentsRow{0pt}{71.\quad Vergl.\ Nr.\ 73, 1797 Juli 27}{.\,.\,.\,.\,.\,.\,.\,.\,.\,.\,.\,.\,.\,.\,.\,.\,.\,.\,.\,.\,.\,.\,.}{—}{521}
\ContentsRow{0pt}{72.\quad Möglichkeit der Ebene bewiesen, 1797 Juli 28}{.\,.\,.\,.\,.\,.\,.\,.\,.\,.\,.\,.\,.\,.\,.\,.}{—}{521}
\ContentsRow{0pt}{73.\quad Einheitswurzeln von zusammengesetztem Grade, 1797 August 1}{.\,.\,.\,.\,.\,.\,.\,.\,.\,.}{—}{521}
\ContentsRow{0pt}{74.\quad Perioden von Einheitswurzeln, 1797 August}{.\,.\,.\,.\,.\,.\,.\,.\,.\,.\,.\,.\,.\,.\,.\,.\,.}{—}{522}
\ContentsRow{0pt}{75.—79.\quad Über Kongruenzen, 1797 August 26, 30, 31, September 4, 9}{.\,.\,.\,.\,.\,.\,.\,.\,.}{—}{522,523}
\ContentsRow{0pt}{80.\quad Erster Beweis der Wurzelexistenz, 1797 Oktober}{.\,.\,.\,.\,.\,.\,.\,.\,.\,.\,.\,.\,.\,.\,.}{—}{523}
\ContentsRow{0pt}{81.\quad Beweis des \textsc{Pythagore}ischen Satzes, 1797 Oktober 16}{.\,.\,.\,.\,.\,.\,.\,.\,.\,.\,.\,.\,.}{—}{524}
\ContentsRow{0pt}{82, 83.\quad Asymptotische Reihenansätze, 1797 Oktober 16, 1798 April}{.\,.\,.\,.\,.\,.\,.\,.\,.\,.}{—}{525}
\ContentsRow{0pt}{84.\quad Klassen binärer quadratischer Formen, 1798 April}{.\,.\,.\,.\,.\,.\,.\,.\,.\,.\,.\,.\,.\,.\,.}{—}{526}
\ContentsRow{0pt}{85.\quad Zusammensetzung von Kräften, 1798 Mai}{.\,.\,.\,.\,.\,.\,.\,.\,.\,.\,.\,.\,.\,.\,.\,.\,.\,.}{—}{529}
\ContentsRow{0pt}{86.\quad Verallgemeinerung des Umkehrungssatzes von \textsc{Lagrange}, 1798 Mai}{.\,.\,.\,.\,.\,.\,.\,.}{—}{530}
\ContentsRow{0pt}{87.\quad Eine besondere Reihe und über allgemeine trigonometrische Reihen, 1798 Juni}{.\,.\,.\,.}{—}{530}
\ContentsRow{0pt}{88.\quad Wahrscheinlichkeitsrechnung, 1798 Juni 17}{.\,.\,.\,.\,.\,.\,.\,.\,.\,.\,.\,.\,.\,.\,.\,.\,.}{—}{533}
\ContentsRow{0pt}{89.\quad Elimination, 1798 Juni}{.\,.\,.\,.\,.\,.\,.\,.\,.\,.\,.\,.\,.\,.\,.\,.\,.\,.\,.\,.\,.\,.\,.\,.\,.}{—}{533}
\ContentsRow{0pt}{90.\quad Anziehung der Kugel, 1798 Juni oder Juli}{.\,.\,.\,.\,.\,.\,.\,.\,.\,.\,.\,.\,.\,.\,.\,.\,.}{—}{533}
\ContentsRow{0pt}{91a, b.\quad Lemniskatische Funktionen, 1798 Juli}{.\,.\,.\,.\,.\,.\,.\,.\,.\,.\,.\,.\,.\,.\,.\,.\,.\,.}{—}{534}
\ContentsRow{0pt}{92.\quad Neues über die Lemniskate, 1798 Juli}{.\,.\,.\,.\,.\,.\,.\,.\,.\,.\,.\,.\,.\,.\,.\,.\,.\,.\,.}{—}{535}
\ContentsRow{0pt}{93.\quad Das ballistische Problem, 1798 Juli}{.\,,\,.\,.\,.\,.\,.\,.\,.\,.\,.\,.\,.\,.\,.\,.\,.\,.\,.\,.}{—}{536}
\ContentsRow{0pt}{94.\quad Theorie der Kometen, 1798 Juli}{.\,.\,.\,.\,.\,.\,.\,.\,.\,.\,.\,.\,.\,.\,.\,.\,.\,.\,.\,.\,.}{—}{536}
\ContentsRow{0pt}{95.\quad Ein neues Feld der Analysis, 1798 Oktober}{.\,.\,.\,.\,.\,.\,.\,.\,.\,.\,.\,.\,.\,.\,.\,.\,.}{—}{536}
\ContentsRow{0pt}{96.\quad Höhere Formen, 1799 Februar 14}{.\,.\,.\,.\,.\,.\,.\,.\,.\,.\,.\,.\,.\,.\,.\,.\,.\,.\,.\,.}{—}{539}
\ContentsRow{0pt}{97.\quad Parallaxe, 1799 April 8}{.\,.\,.\,.\,.\,.\,.\,.\,.\,.\,.\,.\,.\,.\,.\,.\,.\,.\,.\,.\,.\,.\,.\,.}{—}{539}
\ContentsRow{0pt}{98.\quad Numerische Übereinstimmung von $M(\sqrt{2},1)$ und $\dfrac{\pi}{\varpi}$, 1799 Mai 30}{.\,.\,.\,.\,.\,.\,.\,.\,.}{—}{542}
\ContentsRow{0pt}{99.\quad Grundlagen der Geometrie, 1799 September}{.\,.\,.\,.\,.\,.\,.\,.\,.\,.\,.\,.\,.\,.\,.\,.\,.}{—}{543}
\ContentsRow{0pt}{100.\quad Neues über das agM, 1799 November}{.\,.\,.\,.\,.\,.\,.\,.\,.\,.\,.\,.\,.\,.\,.\,.\,.\,.\,.}{—}{544}
\ContentsRow{0pt}{101, 102.\quad Integraleigenschaften des agM, 1799 Dezember 14, 23}{.\,.\,.\,.\,.\,.\,.\,.\,.\,.\,.}{—}{544}
\ContentsRow{0pt}{103.\quad Ternäre Formen, 1800 Februar 13}{.\,.\,.\,.\,.\,.\,.\,.\,.\,.\,.\,.\,.\,.\,.\,.\,.\,.\,.\,.\,.}{—}{545}
\ContentsRow{0pt}{104.\quad Konvergenz einer trigonometrischen Reihe, 1800 April 27}{.\,.\,.\,.\,.\,.\,.\,.\,.\,.\,.\,.}{—}{545}
\ContentsRow{0pt}{105.\quad Mehrdeutigkeit des elliptischen Integrales, 1800 Mai 6}{.\,.\,.\,.\,.\,.\,.\,.\,.\,.\,.\,.\,.}{—}{546}
\ContentsRow{0pt}{106.\quad Bedeutende Ergänzung dieser Theorie, 1800 Mai 22}{.\,.\,.\,.\,.\,.\,.\,.\,.\,.\,.\,.\,.\,.}{—}{547}
\ContentsRow{0pt}{107.\quad Osterformel, 1800 Mai 16}{.\,.\,.\,.\,.\,.\,.\,.\,.\,.\,.\,.\,.\,.\,.\,.\,.\,.\,.\,.\,.\,.\,.\,.}{—}{547}
\Needspace{4\baselineskip}
\ContentsRow{0pt}{108.\quad Der sinus lemniscaticus universalissime acceptus, 1800 Ende Mai bis Juni 2, 3}{.\,.\,.\,.}{—}{548}
\par\noindent X{\small 1}.\hfill 74\par
\endgroup
